\documentclass[11pt,a4paper]{article}
\usepackage{amsmath,amssymb}
\usepackage[margin=2.5cm]{geometry}
\usepackage{array}
\usepackage{bm}
\usepackage{graphicx}
\usepackage{booktabs}
\usepackage{xcolor}
\usepackage{hyperref}
\hypersetup{colorlinks=true, linkcolor=blue!50!black, urlcolor=blue!50!black}

\title{Flow-Conditioned Mamba-3 for 3D Point Tracking}
\author{Masahiro Ogawa}
\date{\today}

\begin{document}
\maketitle

\begin{abstract}
We present a small learned model for 3D point tracking that conditions
on optical-flow vectors and metric-depth values as direct input
features rather than raw image patches. A frozen SEA-RAFT flow model
provides dense 2D correspondence between consecutive frames; a frozen
DA3-Large model provides absolute metric depth. These two sources of
information determine per-point 3D positions up to estimation error.
A small Mamba-3 state-space module then refines the initial 2D
positions by processing the flow-and-depth observation sequence
causally over time. The learned component has approximately 0.5M
parameters compared with 22M in the prior image-encoder-based variant
(v31). We evaluate on TAPVid-3D (minival, three subsets) and compare
against the training-free SEA-RAFT baseline (v31 replacement) and
published BootsTAPIR and SpatialTracker numbers.
\end{abstract}

\tableofcontents

%% ================================================================
\section{Problem: 3D Point Tracking}
%% ================================================================

Given a video $\bm{I} = (I_1, \dots, I_F)$ with $F$ frames, a set of
$N$ query points $\{(x_n, y_n, a_n)\}_{n=1}^{N}$, and camera intrinsic
matrix $\bm{K}$, the goal is to output for each query $n$ and each
frame $t$ a 3D camera-frame position
$\hat{\bm{p}}_{t,n} \in \mathbb{R}^3$ and a visibility score
$\hat{v}_{t,n} \in [0, 1]$.

We use the TAPVid-3D protocol: queries are specified in the original
image's pixel coordinates, ground-truth positions are in the camera
frame of the first frame, and evaluation uses 3D Average Jaccard
(3D-AJ), Average Percentage of correct 3D points within a threshold
(APD3D), and Occlusion Accuracy (OA)~\cite{tapvid3d}.

%% ================================================================
\section{Observations: Optical Flow and Metric Depth}
%% ================================================================

\subsection{SEA-RAFT Optical Flow}

We run a frozen pretrained SEA-RAFT model~\cite{searaft}
(checkpoint: \texttt{Tartan-C-T-TSKH-spring540x960-M}) on consecutive
frame pairs to obtain dense forward flow fields
\begin{equation}
  \bm{f}_t \in \mathbb{R}^{H \times W \times 2}, \qquad t = 1, \dots, F-1,
\end{equation}
where $\bm{f}_t(u, v)$ is the pixel displacement from $(u, v)$ in
frame $t$ to the corresponding point in frame $t+1$.

\subsection{Baseline 2D Track: Flow Chaining}

Starting from the query position $(x_n, y_n)$ at anchor frame $a_n$,
we propagate forward and backward by bilinearly sampling the flow:
\begin{align}
  \bm{q}_{t+1, n}^{(0)} &= \bm{q}_{t,n}^{(0)} +
    \bm{f}_t\bigl(\bm{q}_{t,n}^{(0)}\bigr), \quad t \geq a_n, \\
  \bm{q}_{t-1, n}^{(0)} &= \bm{q}_{t,n}^{(0)} +
    \bm{g}_{t-1}\bigl(\bm{q}_{t,n}^{(0)}\bigr), \quad t \leq a_n,
\end{align}
where $\bm{g}_t$ is the backward flow (frame $t+1 \to t$). Forward-backward
consistency determines a binary visibility mask $v_{t,n}^{(0)} \in \{0, 1\}$.
We refer to this as the \emph{SEA-RAFT baseline} (v-SEARAFT in
Table~\ref{tab:results}).

\subsection{DA3 Metric Depth}

A frozen DA3-Large depth model~\cite{da3} has been applied
offline to all clips, producing per-frame depth maps
$D_t \in \mathbb{R}^{H_d \times W_d}$ stored as quantized
\texttt{uint16} arrays under \texttt{$\sim$/data/tapvid3d\_da3/}.
We dequantize on load:
\begin{equation}
  D_t = d_{\min} + \frac{q_t}{65535}\,(d_{\max} - d_{\min}).
\end{equation}
Sampling depth at a 2D position $\bm{q}_{t,n}$ via bilinear interpolation
gives a per-point scalar $z_{t,n} = D_t(\bm{q}_{t,n})$.

\subsection{Unprojection}

The initial 3D position estimate follows directly from the pinhole
camera model:
\begin{equation}\label{eq:unproject}
  \bm{p}_{t,n}^{(0)} =
  \begin{pmatrix}
    (u_{t,n} - c_x) / f_x \\
    (v_{t,n} - c_y) / f_y \\
    1
  \end{pmatrix}
  z_{t,n},
  \qquad
  \bm{K} = \begin{pmatrix} f_x & 0 & c_x \\ 0 & f_y & c_y \\ 0 & 0 & 1 \end{pmatrix}.
\end{equation}
This is already a strong baseline: SEA-RAFT supplies plausible $u, v$,
DA3-Large supplies absolute scale $z$. Table~\ref{tab:results} shows
this achieves mean 3D-AJ of 5.2\% on TAPVid-3D minival,
an 11$\times$ improvement over the preceding learned variant (v31).

%% ================================================================
\section{Model: Flow-Conditioned Mamba-3 Tracker (v32)}
%% ================================================================

\subsection{Motivation}

The SEA-RAFT baseline still suffers from cumulative drift on long clips
(notably drivetrack with large fast camera motion: 1.6\% vs.\ a
published baseline of 5.1\%). The root cause is that naive chaining
propagates errors monotonically: if the flow at step $t$ is off by
$\epsilon$, the error at step $t + k$ is at least $k\epsilon$ without
any correction.

A Mamba-3 SSM can suppress drift by maintaining a recurrent state that
aggregates evidence from past frames. Crucially, the flow-and-depth
observations at each frame are already informative enough that no
image encoder is needed: the model is trained to learn a temporal
denoising prior directly from the observation sequence.

\subsection{Input Features}

For each track $n$ and frame $t$ we construct a six-dimensional
observation vector:
\begin{equation}\label{eq:phi}
  \bm{\phi}_{t,n} =
  \Bigl[
    \underbrace{\bm{q}_{t,n}^{(0)} / S}_{\text{normalised 2D pos.} \in \mathbb{R}^2},\;\;
    \underbrace{\bm{f}_t(\bm{q}_{t,n}^{(0)}) / S}_{\text{flow vector} \in \mathbb{R}^2},\;\;
    \underbrace{z_{t,n} / z_{\text{ref}}}_{\text{depth} \in \mathbb{R}^1},\;\;
    \underbrace{v_{t,n}^{(0)}}_{\text{vis flag} \in \mathbb{R}^1}
  \Bigr] \in \mathbb{R}^6,
\end{equation}
where $S$ is the image side length (896), $z_{\text{ref}}$ is the
median anchor depth over the batch, and $v_{t,n}^{(0)} \in \{0,1\}$
is the forward-backward consistency mask.

\subsection{Architecture}

\paragraph{Embedding.}
Each observation is mapped to a $D$-dimensional feature:
\begin{equation}
  \bm{e}_{t,n} = \mathrm{MLP}_\mathrm{embed}(\bm{\phi}_{t,n}) \in \mathbb{R}^D,
  \qquad D = 128.
\end{equation}

\paragraph{Causal Temporal Mamba-3.}
We process each track independently as a length-$F$ sequence using $L$
stacked Mamba-3 cross-attention layers in \emph{self-attention mode}
(query and key-value come from the same sequence). Let
$\bm{E}^{(0)}_n = (\bm{e}_{1,n}, \dots, \bm{e}_{F,n}) \in \mathbb{R}^{F \times D}$;
then for $\ell = 1, \dots, L$:
\begin{equation}\label{eq:mamba}
  \bm{E}^{(\ell)}_n = \mathrm{LN}\Bigl(
    \bm{E}^{(\ell-1)}_n + \mathrm{Mamba3}\bigl(
      \mathrm{LN}(\bm{E}^{(\ell-1)}_n),\;
      \mathrm{LN}(\bm{E}^{(\ell-1)}_n)
    \bigr)
  \Bigr),
\end{equation}
where Mamba3$(\bm{Q}, \bm{KV})$ is the variant-B cross-attention from
\texttt{doc/attention/mamba3\_attention.tex}, \S 9:
\begin{equation}
  \bm{Y} = \bigl(\bm{L}_\mathrm{cross} \odot \bm{C}^q \bm{B}^{kv\top}\bigr)\,\bm{V}^{kv},
  \qquad
  L_\mathrm{cross}[i,j] = \gamma_j \prod_{k=j+1}^{F} \alpha_k \cdot \mathbf{1}[j \leq i].
\end{equation}
The causal mask $\mathbf{1}[j \leq i]$ ensures frame $t$ attends only
to past frames, giving the model a left-to-right inductive bias.
All $N$ tracks are processed in parallel by reshaping to $(B \cdot N, F, D)$.

\paragraph{Output Heads.}
After $L$ layers a final LayerNorm produces the refined feature
$\bm{h}_{t,n} \in \mathbb{R}^D$. Two small MLPs read out:
\begin{align}
  \Delta\bm{q}_{t,n} &= \mathrm{MLP}_{uv}(\bm{h}_{t,n}) \in \mathbb{R}^2,
  & \text{(zero-init final layer)} \\
  \hat{v}_{t,n}^{\text{logit}} &= \mathrm{MLP}_{vis}(\bm{h}_{t,n}) \in \mathbb{R}.
\end{align}
The predicted 2D position is
\begin{equation}
  \hat{\bm{q}}_{t,n} = \bm{q}_{a_n,n}^{(0)} + \Delta\bm{q}_{t,n},
\end{equation}
and the final 3D position follows from~\eqref{eq:unproject}
evaluated at $\hat{\bm{q}}_{t,n}$ with depth $D_t(\hat{\bm{q}}_{t,n})$.

The zero-initialisation of the final $\mathrm{MLP}_{uv}$ layer ensures
$\hat{\bm{q}}_{t,n} = \bm{q}_{a_n,n}^{(0)}$ at the start of training,
so the initial loss is bounded by the SEA-RAFT baseline loss rather
than a random-prediction loss.

\subsection{Parameter Count}

\begin{center}
\begin{tabular}{lrr}
\toprule
Component & Parameters & Notes \\
\midrule
$\mathrm{MLP}_\mathrm{embed}$ (6$\to$128)   & $\approx$ 17k & two linear layers \\
Mamba-3 layers $\times 2$  & $\approx$ 430k & dim=128, state=64, heads=4 \\
$\mathrm{MLP}_{uv}$, $\mathrm{MLP}_{vis}$ & $\approx$ 33k & \\
LayerNorms             & $\approx$ 2k & \\
\midrule
\textbf{Total trainable} & $\approx$ \textbf{480k} & vs.\ 22M in v31 \\
\bottomrule
\end{tabular}
\end{center}

%% ================================================================
\section{Training}
%% ================================================================

\subsection{Loss}

We reuse the v31 loss function (TrackingLossV31):
\begin{equation}
  \mathcal{L} = w_{3D}\,\mathcal{L}_{3D} + w_{2D}\,\mathcal{L}_{2D} + w_\text{vis}\,\mathcal{L}_\text{vis},
\end{equation}
with $w_{3D} = w_{2D} = w_\text{vis} = 1$.

The 3D term normalises by the median anchor depth $\bar{z}$ over the
batch (a fixed reference to make the loss scale-invariant):
\begin{equation}
  \mathcal{L}_{3D} = \frac{1}{|\mathcal{V}|}
  \sum_{(t,n) \in \mathcal{V}}
  \frac{\|\hat{\bm{p}}_{t,n} - \bm{p}_{t,n}^*\|_1}{\bar{z}},
\end{equation}
where $\mathcal{V}$ is the set of visible (frame, track) pairs.
The 2D term measures pixel residual normalised by image size $S$:
\begin{equation}
  \mathcal{L}_{2D} = \frac{1}{|\mathcal{V}|}
  \sum_{(t,n) \in \mathcal{V}}
  \frac{\|\hat{\bm{q}}_{t,n} - \bm{q}_{t,n}^*\|_1}{S},
\end{equation}
where $\bm{q}_{t,n}^* = \pi(\bm{p}_{t,n}^*, \bm{K})$ is the GT 2D
projection. The visibility term is binary cross-entropy.

\subsection{Data and Optimiser}

Training uses the TAPVid-3D official split (FULL\_EVAL as training set,
MINIVAL held out for evaluation). We sample 8-frame windows from each
clip during training. Optimiser: AdamW, $\eta = 3 \times 10^{-4}$,
weight decay $0.05$, gradient clip $1.0$. Schedule: warmup 500 steps,
cosine decay for 2000 steps before linear tail to 20k total.
SEA-RAFT and DA3 parameters are frozen throughout.

%% ================================================================
\section{Results}
\label{sec:results}
%% ================================================================

\begin{table}[h]
\centering
\caption{TAPVid-3D \textbf{minival} 3D-AJ (\%) — higher is better.
Baselines marked * are paper numbers from~\protect\cite{tapvid3d}.
Full-eval SoTA is on a different (larger) test set.}
\label{tab:results}
\begin{tabular}{lrrrr}
\toprule
Method & pstudio & drivetrack & adt & mean \\
\midrule
BootsTAPIR + ZoeDepth*   & 10.2 & 5.1  & 8.6  & 8.0 \\
BootsTAPIR + COLMAP*     & 6.2  & 9.3  & 7.3  & 7.6 \\
SpatialTracker*          & 9.8  & 5.8  & 9.2  & 8.3 \\
\midrule
v31 (Mamba-3+DA3, learned)    & 0.5  & 0.2  & 0.7  & 0.5 \\
SEA-RAFT + DA3 (training-free) & 7.5  & 1.6  & 6.6  & 5.2 \\
\textbf{v32 (flow-cond.\ Mamba-3)} &
  \textbf{5.7}  & \textbf{0.4}  & \textbf{4.2}  & \textbf{3.4} \\
\textbf{v33 (depth-refine Mamba-3)} &
  \textbf{2.7}  & \textbf{0.8}  & \textbf{6.1}  & \textbf{3.2} \\
\midrule
\textit{CoTracker3-online (full-eval)} & 21.3 & 10.7 & 14.6 & 15.5 \\
\textit{DELTA (full-eval)}             & 24.4 & 13.0 & 17.6 & 18.3 \\
\bottomrule
\end{tabular}
\end{table}

\noindent \textbf{Result (20k steps, $\sim$9.3\,h).} Contrary to the
hypothesis, the learned refinement \emph{degrades} 3D-AJ on every subset
relative to the training-free SEA-RAFT+DA3 baseline (mean 3.4\% vs.\ 5.2\%).
It does avoid the v31 zero-motion collapse — \texttt{uv\_head} gradients stay
live throughout and the eval motion ratios match SEA-RAFT to within 0.1\%
(100.6/106.9/117.6\%). That near-identical motion magnitude is the diagnosis:
the residual $\Delta uv$ is tiny in aggregate yet consistently adds small
systematic error, most plausibly by nudging $uv$ across DA3 depth
discontinuities (object boundaries) so points flip out of the AJ
within-threshold band. Because SEA-RAFT's $uv$ is already near the ground-truth
2D track, the 2D loss term has almost nothing to correct and training mostly
perturbs a near-optimal input. A residual position delta on an already-good
flow estimate is the wrong lever; promising alternatives are a learned
\emph{consistency gate} (refine only where forward–backward consistency is low)
or a depth-edge-aware 3D loss rather than an unconditional $\Delta uv$.

\medskip
\noindent \textbf{v33: depth-along-ray refinement (2D frozen).} Following the
v32 lesson, the v33 SSM (\S6) freezes \texttt{uv} and the FB-consistency
visibility and refines \emph{only} per-track depth along each pixel ray
($z = z_{\text{raw}}\,e^{\Delta\log z}$), the unique 3D degree of freedom that
leaves the 2D projection invariant. The structural constraint held perfectly:
eval motion ratios ($100.58/106.88/117.53\%$) and occlusion accuracy
($\overline{\text{OA}}=0.728$) are byte-identical to SEA-RAFT. Yet 3D-AJ again
\emph{regresses} (mean $3.2\%$ vs.\ $5.2\%$, worst on pstudio $7.5\!\to\!2.7$).
The diagnosis is a train/eval gap compounded with loss–metric misalignment: the
validation 3D L1 loss falls cleanly ($0.21\!\to\!0.088$ over training) while
minival 3D-AJ drops, because (i) 3D-AJ applies global \emph{median scaling}
before thresholding and therefore scores \emph{relative} depth structure, which
the L1 objective distorts while lowering absolute error, and (ii) the per-point
depth corrections overfit the train depth distribution (pstudio ${\sim}3$\,m,
drivetrack ${\sim}20$\,m, adt ${\sim}1$\,m). Both learned refinements tried ---
2D ($v32$) and depth ($v33$) --- underperform the training-free baseline, which
remains the strongest tracker here. A learned refiner would need a
scale-invariant / median-matched 3D loss and corrections gated to
low-confidence points; otherwise refinement on top of SEA-RAFT+DA3 is not the
right lever.

%% ================================================================
\begin{thebibliography}{9}
\bibitem{tapvid3d}
  Doersch et~al.,
  \textit{TAPVid-3D: A Benchmark for Tracking Any 3D Point},
  NeurIPS 2024.

\bibitem{searaft}
  Wang et~al.,
  \textit{SEA-RAFT: Simple, Efficient, Accurate RAFT for Optical Flow},
  ECCV 2024 (Oral).

\bibitem{da3}
  Ke et~al.,
  \textit{Depth Anything V3},
  arXiv 2025.
\end{thebibliography}

\end{document}
